\begin{note}
	The exact value of $R(k,l)$ remains open. What we know, for now, is that $R(4,4) =18, R(4,5) = 25$. Even $R(5,5)$ is still open.
\end{note}

\begin{proof}
	We prove by induction on $k+l$.

	Base case: when $k=1$ or $l=1$, for all $n\ge 1$, either $\beta(G)\ge k$ or $\omega(G)\ge l$ holds.

	Inductive case: assume the theorem holds for all $k'+l' = k+l-1$.
	Fix a blue-orange coloring on $K_n$. To prove $n\ge \binom{k+l-2}{k-1}$, it suffices to prove the case for $n = \binom{k+l-2}{k-1}$.
	Note that
	\[ \binom{k+l-2}{k-1} = \binom{k+l-3}{k-2} + \binom{k+l-3}{k-1}. \]
	Fix a vertex $v\in K_n$. Let $O = \{u\in K_n: (u,v) \text{ is orange}\}$ and
	$B = \{u\in K_n: (u,v) \text{ is blue}\}$. Then $|O| + |B| = n-1 = \binom{k+l-2}{k-1}-1 = \binom{k+l-3}{k-2} + \binom{k+l-3}{k-1}-1$.

	We claim, either $|O|\ge \binom{k+l-3}{k-1}$ or $|B|\ge \binom{k+l-3}{k-2}$. Otherwise, we have
	\[ |O| \le \binom{k+l-3}{k-1} - 1 \text{ and } |B| \le \binom{k+l-3}{k-2} -1 \implies |O| + |B| \le n-2 \implies \contra. \]

	\begin{itemize}
		\item If $|O|\ge \binom{k+l-3}{k-1}$, then by the inductive hypothesis, either $\beta(O)\ge k$ or $\omega(O)\ge l-1$. If $\beta(O)\ge k$, then we have $\beta(G)\ge k$. Otherwise, if $\omega(O)\ge l-1$, then together with $v$, we have $\omega(G)\ge l$.
		\item If $|B|\ge \binom{k+l-3}{k-2}$, then by the inductive hypothesis, either $\beta(B)\ge k-1$ or $\omega(B)\ge l$. If $\beta(B)\ge k-1$, then together with $v$, we have $\beta(G)\ge k$. Otherwise, if $\omega(B)\ge l$, then we have $\omega(G)\ge l$.
	\end{itemize}
	In all cases, we have either $\beta(G)\ge k$ or $\omega(G)\ge l$. Therefore, $R(k,l) \le n = \binom{k+l-2}{k-1}$.
\end{proof}

The same strategy can be used to prove the following (we omit here):
\begin{theorem}
	\[ R(k,l) \le R(k,l-1) + R(k-1,l).\]
\end{theorem}
